%% AlpacaVision AI — Cover Letter for MDPI Animals
%% Corresponding author: Leonid Aleman-Gonzales (advisor)
%% Portal: susy.mdpi.com

\documentclass[12pt, letterpaper]{article}
\usepackage[margin=2.5cm]{geometry}
\usepackage{hyperref}
\usepackage{parskip}
\usepackage{microtype}

\hypersetup{colorlinks=true, urlcolor=blue, linkcolor=blue}

\begin{document}

\begin{flushright}
  Leonid Alem\'an-Gonzales (Corresponding Author)\\
  Escuela Profesional de Ingenier\'ia Estad\'istica e Inform\'atica\\
  Universidad Nacional del Altiplano (UNAP)\\
  Puno, Per\'u\\
  \href{mailto:laleman@unap.edu.pe}{laleman@unap.edu.pe}\\[6pt]
  \today
\end{flushright}

\bigskip

The Editor-in-Chief\\
\textit{Animals}\\
MDPI AG, Basel, Switzerland\\
\href{https://www.mdpi.com/journal/animals}{https://www.mdpi.com/journal/animals}

\bigskip

Dear Editor,

We are pleased to submit our manuscript entitled:

\begin{center}
\textbf{Automatic Computer-Vision Detection of Morphological Anomalies in \\
Altiplano Alpacas: A Curated Dataset, Compact Detector, and \\
Ocular Classification Feasibility Study}
\end{center}

for consideration for publication in \textit{Animals}.

\textbf{Scientific contribution.}
Peru holds 87\% of the world's alpaca (\textit{Vicugna pacos}) population, yet no
publicly available computer-vision dataset or automated diagnostic system exists for
this economically vital species. This paper presents AlpacaVision AI, comprising a
curated open dataset, a compact and deployable body detector, a data-integrity
methodology, and a rigorous, honest feasibility study of automatic ocular anomaly
classification. Our compact YOLOv11n detector reaches mAP@0.5 = 0.860 on a
leakage-free 308-image held-out test set; a larger YOLOv11s reaches only 0.863,
confirming the compact model captures essentially all available signal.

\textbf{Key contributions.}
\begin{enumerate}
  \item \textbf{Open, curated dataset:} 2,051 unique labelled alpaca images (YOLO
        format), consolidated from nine public sources after removing 1,037 exact
        MD5 duplicates, with CLAHE preprocessing for Andean high-UV conditions,
        released under CC~BY~4.0.
  \item \textbf{Compact detector:} YOLOv11n (2.58M parameters) achieving
        mAP@0.5 = 0.860 (Precision = 0.913, Recall = 0.731) on a strictly
        leakage-free held-out test set, with ONNX export for field deployment.
  \item \textbf{Data-integrity methodology:} MD5 deduplication, group-aware
        train/validation/test splitting, and label-conflict resolution, developed
        after an audit revealed systematic data leakage that had inflated all
        initially reported metrics --- a transparent, reproducible protocol we
        release with the code.
  \item \textbf{Honest feasibility study (negative result):} under a leakage-free
        protocol the EfficientNet-B2 ocular classifier performs at chance
        (AUC-ROC = 0.506, 95\% CI [0.368, 0.645]). Because the identical pipeline
        reaches F1 = 0.957 on real human-fundus data, we attribute the failure to
        the absence of expert veterinary labels rather than to the model, and
        conclude that vision-language auto-labelling alone is insufficient for
        fine-grained camelid ocular pathology.
  \item \textbf{Deployed system:} an end-to-end Flask web application with role-based
        access and automated Spanish-language veterinary report generation.
\end{enumerate}

\textbf{Why this matters.}
Beyond the first open alpaca detection benchmark, the manuscript directly addresses
\textit{Animals}' core themes of animal health, welfare, and the practical management of
livestock: it provides a deployable tool to support early, objective detection of
morphological anomalies in a species whose remote high-altitude husbandry offers minimal
access to veterinary care, and it discusses concrete pathways to integration in field
health campaigns and herd-welfare monitoring. Its candid treatment of data leakage and its
well-documented negative result are of additional broad value, setting realistic
expectations for practitioners tempted to bootstrap medical-grade classifiers from
large-language-model annotations and offering a reusable methodology for leakage-free
evaluation in veterinary auto-labelling pipelines. The work thus sits squarely within the
scope of \textit{Animals} at the intersection of precision livestock farming, animal
welfare, and AI-assisted veterinary diagnostics, with strong applied relevance for
indigenous Andean herding communities.

\textbf{Originality statement.}
This manuscript has not been published and is not under consideration elsewhere.
The dataset and code are openly released (CC~BY~4.0 and MIT, respectively).
All co-authors have approved the final manuscript and declare no competing interests.
The article-processing charge is supported by the Universidad Nacional del Altiplano
(UNAP), Puno, Peru.

\textbf{Suggested reviewers.}
We suggest reviewers with expertise in:
(1)~computer vision applied to livestock health and behaviour monitoring;
(2)~transfer learning and convolutional neural networks for animal disease detection;
(3)~precision livestock farming technology and its adoption in developing-country contexts.

We thank you for considering this submission and look forward to your response.

Sincerely,

\vspace{1.5cm}

Leonid Alem\'an-Gonzales (Advisor / Corresponding Author)\\[4pt]
Richar Andre Vilca-Solorzano\\
Dina Maribel Yana-Yucra\\
Cristian Daniel Ccopa-Acero\\[6pt]
Semillero de Investigaci\'on ``John J. Hopfield -- IIICCD''\\
Universidad Nacional del Altiplano (UNAP), Puno, Per\'u

\end{document}
